\subsection{GDL 随机微结构生成方法与参数设置}

\subsubsection{微结构生成核心框架}
本文基于分层堆叠与各向异性扰动策略，开发了GDL碳纤维随机微结构生成算法，解决传统生成方法中“各向同性毛球效应”“体素硬边界锯齿化”“粘结剂（binder）散点分布”等问题。算法核心流程为：分层划分→纤维几何采样（方向+中心）→纤维体素化（支持重叠/不重叠模式）→binder 精准生成（交点胶团+毛细桥约束）→体素场平滑，最终输出0~1连续体素场（0为孔隙，1为固相），提升微结构真实性与可视化平滑度。

\subsubsection{基础几何与孔隙率参数（表2）}
微结构计算域尺寸设置为 \(I \times J \times K = 155 \times 300 \times 300\) 体素（对应物理尺度 µm 级，厚度轴自动识别为最小维度，即 \(I\) 轴），总孔隙数 \(N=I \times J \times K\)。目标总孔隙率（含纤维与binder）设定为 \(\varepsilon_{\text{target,total}} = 0.8018\)，孔隙率收敛精度 \(\Delta\varepsilon = 0.0005\)，避免生成过程中孔隙率波动过大。

碳纤维几何参数参考实际GDL碳纤维特性，具体设置如表\ref{tab:basic_params} 所示。

\begin{table}[H]
\centering
\small
\begin{adjustbox}{width=\columnwidth}
\begin{tabular}{lll}
\toprule
\textbf{符号/名称} & \textbf{含义} & \textbf{典型取值/范围} \\
\midrule
\(I,J,K\) & 体素域尺寸 & \(155 \times 300 \times 300\) \\
\(r_{\text{fiber}}\) & 纤维半径（体素） & 3.5 \\
\(L_{\text{min}}, L_{\text{max}}\) & 纤维长度范围 & 140.0，320.0 \\
\texttt{fiber\_end\_mode} & 端点处理方式 & \texttt{clipped} \\
\texttt{max\_fibers} & 纤维生成上限 & 50000 根 \\
\(\varepsilon_{\text{target,total}}\) & 目标总孔隙率 & 0.8018 \\
\(\Delta\varepsilon\) & 孔隙率收敛精度 & 0.0005 \\
\bottomrule
\end{tabular}
\end{adjustbox}
\caption{基础几何与孔隙率参数}
\label{tab:basic_params}
\end{table}

{\raggedright
\hspace{2em}纤维长度采用长纤维裁剪模式（\texttt{clipped}），该模式先生成长纤维中心线再裁剪至计算域，使纤维端点主要分布在计算域边界，更贴近真实GDL子区域截取特征。
\par}
% 修正后的表4：自适应列宽 + 限制总宽度 + 自动换行
\subsection{分层与各向异性参数（表3）}
为模拟GDL“分层堆叠”的宏观特征，沿厚度轴（$I$ 轴）划分分层数 $n_{\text{layers}}$，分层厚度 $h_{\text{layer}} = \max(4.0 \times r_{\text{fiber}}, 1.0)$ 体素，最终分层数取 $\max(3, \text{round}(\text{thickness}/h_{\text{layer}}))$（至少3层保证分层效果）。层内纤维中心在厚度轴方向的扰动（jitter）设为 $0.35 \times h_{\text{layer}}$，平衡分层规整性与真实随机性。

各向异性扰动参数控制纤维方向的“乱中有序”，具体设置如表 \ref{tab:anisotropic_params} 所示。
\begin{table}[H]
\centering
\Huge
\begin{adjustbox}{width=\columnwidth}
\begin{tabular}{lll}
\toprule
\textbf{符号/名称} & \textbf{含义} & \textbf{取值/范围} \\
\midrule
$n_{\text{layers}}$ & 分层数 & $\max(3, \text{round}(\text{thickness}/h_{\text{layer}}))$ \\
$\sigma_z$ (jitter) & 层内厚度方向抖动 & $0.35 \times h_{\text{layer}}$ \\
$\sigma_{\alpha}$ & 面内角度离散度 & $55.0^\circ$ \\
$\sigma_{\beta}, \beta_{\max}$ & 离层倾角控制 & $\sigma_{\beta}=4.0^\circ,\ \beta_{\max}=10.0^\circ$ \\
相邻层方向步长 & 层间方向变化步长 & $20^\circ \text{（随机游走）}$ \\
\bottomrule
\end{tabular}
\end{adjustbox}
\caption{分层与各向异性参数}
\label{tab:anisotropic_params}
\end{table}

{\raggedright
\hspace{2em}面内角度扰动标准差\(\sigma_{\alpha}=55.0^\circ\)，该值越大面内方向越随机，当取值为 \(180^\circ\) 时近似各向同性；离层面倾角扰动标准差 \(\sigma_{\beta}=4.0^\circ\)，倾角硬截断阈值 \(\beta_{\text{max}}=10.0^\circ\)，可限制纤维沿厚度轴的分量，保证纤维主要平行于层面。每层平均面内方向采用“随机游走”生成，相邻层平均方向步长 \(20^\circ\)，使层间方向缓慢变化，强化分层视觉特征。
\par}


\subsection{binder 生成核心参数（表4）}
{\raggedright
\hspace{2em}pacbinder 体积分数（相对总体）设定为 \(\phi_{\text{binder,total}} = 1.2\%\)（0.012），参考 GDL 中 binder 占固相体积的典型比例，具体参数如表 \ref{tab:binder_params} 所示。
\par}
\begin{table}[H]
\centering
\small
\begin{adjustbox}{width=\columnwidth}
\begin{tabular}{lll}
\toprule a
\textbf{符号/名称} & \textbf{含义} & \textbf{取值/范围} \\
\midrule
\(\phi_{\text{binder,total}}\) & binder体积分数 & 0.012 \\
\(g_{\text{contact}}\) & 接触间隙 & \(0.35 \times r_{\text{fiber}}\) \\
\(r_{\text{binder,blob}}\) & 初始胶团半径 & \(1.05 \times r_{\text{fiber}}\) \\
\(d_{\text{user}}\) & 毛细桥连线长度 & \(2.65 \times r_{\text{fiber}}\) \\
\(\text{d}t, \text{d}s\) & 采样密度（体素） & 0.65，0.75 \\
约束条件 & 体素范围约束 & 纤维间隙slab区域 \\
\bottomrule
\end{tabular}
\end{adjustbox}
\caption{binder 生成核心参数}
\label{tab:binder_params}
\end{table}

{\raggedright
\hspace{2em}binder 生成时，仅在两根纤维“夹缝区”生成，体素需同时满足到两根纤维中心线距离 \(\leq \text{thr}\)（\(\text{thr}^2\) 为距离平方阈值），且限制在纤维间隙的slab区域内（投影范围 \([-{\text{margin}}, g+{\text{margin}}]\)），避免binder在单纤维表面形成凸起，保证胶桥的凹形圆角特征。
\par}

\subsection{数值计算与优化参数（表5）}
为提升算法效率与结果可复现性，采用多项数值优化策略，具体设置如表 \ref{tab:optimization_params} 所示。
\\
% 修正后的表5：同样适配单栏宽度
\begin{table}[H]
\centering
\begin{tabular}{p{1.8cm}p{2.7cm}p{2.5cm}}
\toprule
策略名称 & 含义 & 设置方式 \\
\midrule
Numba JIT 编译 & 核心几何计算加速 & 加速 point segment dist2 等函数 \\
AABB 分箱预筛 + 微调 & 不重叠模式优化 & 降低高密度碰撞检测开销 \\
Scipy 距离变换 & 体素平滑方式 & 生成窄过渡带（无则降级） \\
全局随机种子 & 可复现性控制 & 42 \\
\bottomrule
\end{tabular}
\caption{数值计算与优化参数}
\label{tab:optimization_params}
\end{table}

采用Numba即时编译（JIT）加速点-线段距离平方（\texttt{point\_segment\_dist2}）、线段-线段最近距离（\texttt{\_segseg\_closest}）等核心几何计算函数，可显著提升体素化效率；不重叠模式下采用轴对齐包围盒（AABB）分箱预筛+厚度轴微调（nudge）策略，能有效降低高密度下纤维碰撞检测的计算开销；体素平滑依赖Scipy的距离变换（\texttt{ndimage}）生成窄过渡带，若无Scipy则降级为无平滑模式，输出连续体素场替代传统0/1/2硬边界，提升Tecplot等值面可视化的光滑度；全局随机种子设为42，保证生成结果的可复现性。

\urlstyle{same}
\subsection{输出设置（表6）}
{\raggedright
\hspace{2em}算法生成4类微结构文件（\texttt{gdl\_layered\-\_opt\-\{opt\}.dat}），覆盖“重叠/不重叠”与“有/无 binder”的组合，具体如表 \ref{tab:output_files} 所示。
\par}

\begin{table}[H]
\centering
\begin{adjustbox}{width=\columnwidth}
\Large
\begin{tabular}{lll}
\toprule
\textbf{文件编号} & \textbf{文件名} & \textbf{组合模式} \\
\midrule
opt=1 & \path{gdl_layered_opt1.dat} & 不重叠+无 binder \\
opt=2 & \path{gdl_layered_opt2.dat} & 可重叠+无 binder \\
opt=3 & \path{gdl_layered_opt3.dat} & 不重叠+有 binder \\
opt=4 & \path{gdl_layered_opt4.dat} & 可重叠+有 binder \\
\bottomrule
\end{tabular}
\end{adjustbox}
\caption{输出文件设置}
\label{tab:output_files}
\end{table}
{\raggedright
\hspace{2em}所有文件均输出0~1连续体素场，可直接导入Tecplot、ParaView等可视化软件进行后处理，或用于孔隙尺度流动/传质模拟。
\par}